% Generated draft: requires mathematical and source review.
\documentclass[11pt]{article}
\usepackage[margin=25mm]{geometry}
\usepackage{fontspec}
\setmainfont{Latin Modern Roman}
\usepackage{amsmath,amssymb,mathtools,unicode-math}
\setmathfont{Latin Modern Math}
\usepackage{enumitem,needspace,xurl,hyperref}
\setlength{\parindent}{0pt}
\setlength{\parskip}{5pt}
\setlength{\emergencystretch}{3em}
\newcommand{\N}{\mathbb{N}}
\newcommand{\Z}{\mathbb{Z}}
\newcommand{\Q}{\mathbb{Q}}
\newcommand{\R}{\mathbb{R}}
\newcommand{\C}{\mathbb{C}}
\newcommand{\F}{\mathbb{F}}
\newcommand{\A}{\mathbb{A}}
\newcommand{\PP}{\mathbb{P}}
\newcommand{\E}{\mathbb{E}}
\newcommand{\eps}{\varepsilon}
\newcommand{\dd}{\,\mathrm{d}}
\begin{document}
\section*{UnsolvedMath ID: 701270}
\textbf{Source problem code:} catalogue entry
\subsection{Definitions and mathematical statement}
A group \(G\) is linear if there exist a field \(K\), a positive integer \(n\), and an injective group homomorphism \(\rho\colon G\hookrightarrow \operatorname{GL}_n(K)\), where \(\operatorname{GL}_n(K)\) is the group of invertible \(n\times n\) matrices over \(K\). A group is simple if it is nontrivial and its only normal subgroups are \(\{1_G\}\) and \(G\). An involution is an element of order exactly \(2\).

Does there exist a group \(G\), a field \(K\), a positive integer \(n\), and an injective homomorphism \(\rho\colon G\hookrightarrow\operatorname{GL}_n(K)\) such that \(G\) is simple,
\[
\exists x,y\in G\quad xy\ne yx,
\]
and
\[
\forall g\in G\quad \bigl(g^2=1_G\Longrightarrow g=1_G\bigr)?
\]
The source specifies neither the base field nor its characteristic; the formulation above uses the usual meaning of linearity, allowing any field and any positive finite matrix dimension. No finiteness assumption on \(G\) is imposed.
\subsection{Short English statement}
Does a nonabelian simple group with no element of order \(2\) admit a faithful representation by invertible matrices of some positive finite dimension over some field?
\subsection{Sources}
\begin{itemize}
\item Local UnsolvedMath registry, permanent ID \(701270\), exact source problem code \texttt{catalogue entry}, title \emph{Kourovka Notebook Problem 16.78}. The supplied registry labels the problem open; this label is not an independently verified mathematical or literature assessment.
\item \emph{The Kourovka Notebook (Unsolved Problems in Group Theory)}, original problem 16.78, 16th issue (2006); attributed to B. Poizat. The supplied record also identifies the 21st issue edition (2026), \texttt{arXiv:1401.0300v48}. The collection description names E. I. Khukhro and V. D. Mazurov as editors. These bibliographic details are supplied attribution and have not been independently checked.
\end{itemize}
\end{document}
